\documentclass[10pt]{article}
\usepackage[T1]{fontenc}
\usepackage[utf8]{inputenc}
\usepackage[french]{babel}
\usepackage{amsmath,amssymb,mathrsfs,mathpazo}
\usepackage{geometry}
\geometry{a4paper,margin=2cm}
\usepackage{graphicx,array,tabularx,tabularray,colortbl,multirow,multicol,booktabs,enumitem,pifont,ulem,fancybox,diagbox,cancel}
\usepackage{pstricks,pstricks-add,pst-plot,pst-tree,pst-node,pst-text,pst-func,pst-math,pst-bspline,pst-3dplot,pst-eucl,pst-solides3d}
\let\pscalculate\undefined
\usepackage{pst-bezier}
\usepackage{tikz,pgfplots,tkz-tab}
\usetikzlibrary{arrows,arrows.meta,patterns,calc,positioning,trees,decorations.markings,shapes.geometric,shapes.misc,shapes.arrows,intersections,decorations.pathmorphing,decorations.pathreplacing,backgrounds,shadows,babel}
\pgfplotsset{compat=1.18}
\usepgfplotslibrary{fillbetween}
\usepackage[np]{numprint}
\usepackage{listings,makecell,pas-tableur,twoopt}
\let\barre\undefined
\usepackage[pstricks]{tablvar}
\definecolor{darkgreen}{RGB}{0,100,0}
\providecommand{\red}{\color{red}}
\providecommand{\blue}{\color{blue}}
\providecommand{\green}{\color{green}}
\providecommand{\black}{\color{black}}
\providecommand{\white}{\color{white}}
\providecommand{\gray}{\color{gray}}
\providecommand{\cyan}{\color{cyan}}
\providecommand{\magenta}{\color{magenta}}
\providecommand{\yellow}{\color{yellow}}
\providecommand{\violet}{\color{violet}}
\providecommand{\orange}{\color{orange}}
\providecommand{\euro}{\texteuro}
\providecommand{\EUR}{\texteuro}
\providecommand{\ieme}{\textsuperscript{e}}
\providecommand{\degres}{\ensuremath{^\circ}}
\providecommand{\degree}{\ensuremath{^\circ}}
\providecommand{\vv}[1]{\overrightarrow{#1}}
\providecommand{\Coord}[1]{#1}
\providecommand{\decosix}{}
\providecommand{\starredbullet}{\ensuremath{\bullet}}
\pagestyle{empty}
\setlength{\parindent}{0pt}
\usepackage{adjustbox}\DeclareUnicodeCharacter{2236}{:}\DeclareUnicodeCharacter{2190}{\ensuremath{\leftarrow}}
\providecommand{\R}{\mathbb{R}}
\providecommand{\N}{\mathbb{N}}
\providecommand{\Z}{\mathbb{Z}}
\providecommand{\Q}{\mathbb{Q}}
\providecommand{\D}{\mathbb{D}}
\providecommand{\C}{\mathbb{C}}
\providecommand{\e}{{\mathrm{e}}}
\providecommand{\vect}[1]{\overrightarrow{#1}}
\providecommand{\vectt}[1]{\overrightarrow{\mathrm{#1}}}
\providecommand{\vvt}[1]{\overrightarrow{\mathrm{#1}}}
\providecommand{\np}[2][]{{#2~#1}}
\providecommand{\vv}[1]{{\overrightarrow{#1}}}
\renewcommand{\d}{{\mathrm{d}}}
\providecommand{\degres}{{^\circ}}
\providecommand{\degre}{{^\circ}}
\providecommand{\blue}{\color{blue}}
\providecommand{\red}{\color{red}}
\providecommand{\enskip}{\quad}
\providecommand{\llbracket}{\lbrack\!\lbrack}
\providecommand{\rrbracket}{\rbrack\!\rbrack}
\providecommand{\euro}{{\text{€}}}
\providecommand{\no}{{\text{n°}}}
\providecommand{\interval}[2]{{[#1;#2]}}
\providecommand{\intervFO}[2]{{[#1;#2[}}
\providecommand{\dotfill}{\ldots}

\providecommand{\ds}{\displaystyle}
\providecommand{\textding}[1]{#1}
\providecommand{\Oij}{\(\mathrm{O};\vec{\imath},\vec{\jmath}\)}
\providecommand{\Oijk}{\(\mathrm{O};\vec{\imath},\vec{\jmath},\vec{k}\)}
\providecommand{\Ouv}{\(\mathrm{O};\vec{u},\vec{v}\)}
\providecommand{\intervFF}[2]{\ensuremath{[#1;#2]}}
\renewcommand{\euro}{\ifmmode\text{\texteuro}\else\texteuro\fi}
\DeclareUnicodeCharacter{00B0}{\ifmmode^{\circ}\else\textdegree\fi}
\DeclareUnicodeCharacter{20AC}{\ifmmode\text{\texteuro}\else\texteuro\fi}
% Énoncés : Éducation nationale. Transcription et corrigés : APMEP, Denis Vergès et contributeurs. Découpage et classement : CoopMaths, Sébastien Lozano, Éric Elter et contributeurs. Adaptation de la présentation : Maths Explorer.
% CC BY-SA : https://coopmaths.fr/www/about/
% https://coopmaths.fr/alea/static/e3c/2021/tex/e3c_2021_01_specimen2_1.tex
\begin{document}
\noindent\textit{APMEP / CoopMaths. CC BY-SA. Présentation adaptée par Maths Explorer.}\par\medskip
\bigskip

\emph{Cet exercice est un questionnaire à choix multiples (QCM). Les cinq questions sont indépendantes. Pour chacune des questions, une seule des quatre réponses est exacte.\\
Le candidat indiquera sur sa copie le numéro de la question et la lettre correspondant à la réponse exacte.\\
Aucune justification n'est demandée. Une réponse exacte rapporte un point, une réponse fausse ou une absence de réponse ne rapporte ni n'enlève aucun point.}

\bigskip

\textbf{Question 1}

\medskip

Pour $x$ pièces produites, le coût de fabrication $C(x)$, en milliers d'euros est donné par

$C(x)= 0,01x^3 - 0,135x^2 + 0,6x +15$, avec $x\in [0~;~30]$.

Pour 2 pièces produites, le coût de fabrication en euros est:

\medskip

{\renewcommand{\arraystretch}{1.5}
\begin{tabularx}{\linewidth}{|X|X|X|X|}
\hline
\textbf{a.}\quad $15,74$ &  \textbf{b.}\quad $157,4$ &  \textbf{c.}\quad $\np{1574}$ &  \textbf{d.}\quad $\np{15740}$\\
\hline
\end{tabularx}}

\bigskip

\textbf{Question 2}

\medskip

Soit $f$ une fonction polynôme du second degré donnée, pour tout nombre réel $x$ par

$f(x)= ax^2+bx+c$, où $a$, $b$, $c$ sont réels. On note $\Delta$ son discriminant. On donne ci-dessous $\mathcal{C}_f$ la courbe représentative de $f$ et on suppose qu'elle admet l'axe des abscisses comme tangente en un de ses points.

\begin{center}
\scalebox{0.7}{
\psset{xunit=1cm, yunit=0.5cm}
\def\xmin {-10.9}   \def\xmax {1.8}
\def\ymin {-15}   \def\ymax {3.2}
\begin{pspicture*}(\xmin,\ymin)(\xmax,\ymax)
\psset{yMaxValue=\ymax,yMinValue=\ymin}
\psgrid[unit=1cm,subgriddiv=1, gridlabels=0, gridcolor=lightgray](-11,-8)(2,3)
\psaxes[labelFontSize=\displaystyle,arrowsize=3pt 2, ticksize=-2pt 2pt, Dy=2]{->}(0,0)(\xmin,\ymin)(\xmax,\ymax)[$x$,-120][$y$,210] 
\def\f{-12 x x 10 add mul 25 add mul 25 div}
\psplot[plotpoints=2000,linewidth=1.25pt]{\xmin}{\xmax}{\f}
\end{pspicture*}
}
\end{center}

On peut affirmer que:

\medskip

{\renewcommand{\arraystretch}{1.5}
\begin{tabularx}{\linewidth}{|X|X|X|X|}
\hline
\textbf{a.}\quad $a<0$ et $\Delta<0$
& \textbf{b.}\quad $a>0$ et $\Delta=0$ 
& \textbf{c.}\quad $a<0$ et $\Delta=0$
& \textbf{d.}\quad $a<0$ et $\Delta>0$\\
\hline
\end{tabularx}}

\bigskip

\textbf{Question 3}

\medskip

$\cos \left (x+\dfrac{\pi}{2}\right)$ est égal à :

\medskip

{
\begin{tabularx}{\linewidth}{|X|X|X|X|}
\hline
\textbf{a.}\quad $\cos(x)-\sin(x)$
& \textbf{b.}\quad $\cos \left (x-\dfrac{\pi}{2}\right )$\rule[-10pt]{0pt}{25pt}
& \textbf{c.}\quad $\sin(x)$
& \textbf{d.}\quad $-\sin(x)$\\
\hline
\end{tabularx}}



\textbf{Question 4}

\medskip

Dans le plan rapporté à un repère orthonormé, on donne les points A\,$(-7~;~4)$ et B\,$(1~;~-2)$.

Le cercle $\Gamma$ de diamètre [AB] admet comme équation dans ce repère :

\medskip

{\renewcommand{\arraystretch}{1.5}
\begin{tabularx}{\linewidth}{|X|X|}
\hline
\textbf{a.}\quad $(x+7)^2+(y-4)^2=100$
& \textbf{b.}\quad $(x+3)^2+(y-1)^2=25$\\
\hline
 \textbf{c.}\quad $(x+3)^2+(y-1)^2=100$
& \textbf{d.}\quad $(x+7)^2+(y-4)^2=25$\\
\hline
\end{tabularx}}

\bigskip

\textbf{Question 5}

\medskip

Dans le plan rapporté à un repère orthonormé, les droites $\mathcal{D}$ et $\mathcal{D}'$ d'équations cartésiennes respectives $3x+2y-1=0$  et $6x+4y+2=0$ sont :

\medskip

{\renewcommand{\arraystretch}{3}
\begin{tabularx}{\linewidth}{|X|X|X|X|}
\hline
\shortstack{\textbf{a.}~ sécantes et non\\~perpendiculaires}
& \shortstack{\textbf{b.}~ confondues\\\phantom{confondues}}
& \shortstack{\textbf{c.}~ strictement\\parallèles}
& \shortstack{\textbf{d.}~ perpendiculaires\\\phantom{confondues}}\\
\hline
\end{tabularx}}

\vspace{0.5cm}

\end{document}
